% !TeX root = ../../main.tex
\subsection{Signed loop moments}
\label{subsec:a:02}

The two flags of a self-loop share one integration variable in
\eqref{eq:s03:integrated-graph-formula}. We retain this correlation by
assigning a weight to their ordered pair of descendant degrees. For
$a,b\in\mathbb Z_{\geq0}$ define
\begin{equation}\label{eq:a:paired-weight}
 w_{ab}=\frac12\int_0^{\alpha}
       \frac{(-t/2)^{a+b}}{(2a+1)!(2b+1)!}\,\dd t
 =\frac{(-1)^{a+b}\pi^{2a+2b+2}}
 {(a+b+1)(2a+1)!(2b+1)!}.
\end{equation}
For a list of $p\geq0$ ordered pairs $(a_i,b_i)$ put
$D=\sum_i(a_i+b_i)$ and let $m_a$ count the endpoints of degree $a$,
for $a=0,1,2$. For a polynomial $\Phi$ in $p,D,m_0,m_1,m_2$ define
the signed linear functional
\begin{equation}\label{eq:a:signed-functional}
 \mathcal L[\Phi]=e^{-S(1)}\sum_{p\geq0}\frac1{p!}
 \sum_{a_1,b_1,\ldots,a_p,b_p\geq0}
       \left(\prod_{i=1}^p w_{a_i b_i}\right)
       \Phi(p,D,m_0,m_1,m_2).
\end{equation}
The empty list has all five statistics zero. The factor $1/2$ in
$w_{ab}$ and $1/p!$ are exactly the root-loop automorphism factors;
there is no additional symmetry factor in this definition.

The following generating function gives all the required moments by
finitely many differentiations. Its variables $v,x,z_0,z_1,z_2$ are
complex moment parameters, independent of the edge variable $t$.

\begin{lmm}[Joint signed moments]\label{lmm:a:joint-moments}
The sums defining $\mathcal L$ are absolutely convergent and
$\mathcal L[1]=1$. More generally,
\begin{align}
 &\mathcal L\left[e^{vp+xD+\sum_{a=0}^2z_a m_a}\right]\notag\\
 &\quad=\exp\left\{\frac{e^v}{2}\int_0^\alpha
 \left[f(e^x t)+\sum_{a=0}^2(e^{z_a}-1)
               \frac{(-e^x t/2)^a}{(2a+1)!}\right]^2\dd t-S(1)\right\}.
 \label{eq:a:joint-egf}
\end{align}
Both sides are entire in the five moment parameters. They define a
signed generating series, with no probability interpretation.
\end{lmm}
\begin{prf}
Expand the square. Selecting the degrees $a,b$ on one edge gives
$w_{ab}e^{x(a+b)}$ multiplied by the two corresponding low-degree
factors. Summing over $p$ labelled pairs with $1/p!$ exponentiates
this one-edge sum and proves the identity. On any compact parameter
set, replacing every coefficient by its absolute value is bounded by
the corresponding factorial series on the finite interval
$[0,\alpha]$. The exponential sum over $p$ then converges normally,
also after any fixed number of derivatives. At zero parameters the
one-edge integral is $\tfrac12\int_0^\alpha f(t)^2\dd t=S(1)$,
which proves the normalization.
\end{prf}

Here are explicit finite rules for differentiating
\eqref{eq:a:joint-egf}. They are also convenient checks on endpoint
correlations. For $a=0,1,2$ and integer $r\geq0$ put
\begin{equation}\label{eq:a:side-moments}
 M_{a,r}=\left.(z\partial_z)^r
 \left\{\frac12\int_0^\alpha
       \frac{(-zt/2)^a}{(2a+1)!}f(zt)\,\dd t\right\}\right|_{z=1}.
\end{equation}
The expression in braces is
\[
 \frac{2(-1)^a}{(2a+1)!z}
              \int_0^{\pi\sqrt z}y^{2a}\sin y\,\dd y.
\]
The integral identity first holds for positive real $z$ and then by
analytic continuation near $1$. Integration by parts and differentiation
give
\begin{equation}\label{eq:a:side-table}
\begin{array}{@{}c|c@{\;}c@{\;}c@{}}
 a&M_{a,0}&M_{a,1}&M_{a,2}\\\hline
0&4&-4&4-\pi^2/2\\
1&(4-\pi^2)/3&(\pi^2-4)/3&(\pi^4-4\pi^2+16)/12\\
2&(\pi^4-12\pi^2+48)/60&-(\pi^4-12\pi^2+48)/60
 &-(\pi^2-4)(\pi^4+48)/240
\end{array}
\end{equation}

For a nonzero nonnegative integer multi-index
$\mathbf r=(r_p,r_D,r_0,r_1,r_2)$ let
$\chi_{\mathbf r}$ be the derivative at zero of the logarithm of
\eqref{eq:a:joint-egf}, of these respective orders. Equivalently it is
the one-loop weighted sum of
$(a+b)^{r_D}\prod_{c=0}^2(\mathbf1_{a=c}+\mathbf1_{b=c})^{r_c}$;
any number of $p$ derivatives supplies a factor one. With no low-count
derivative the one-loop degree moments, for $r_D=0,\ldots,4$, are
\begin{equation}\label{eq:a:degree-cumulants}
 S(1),\quad -S(1),\quad S(1),\quad
 \frac{\pi^2}{2}-S(1),\quad S(1)-\frac{5\pi^2}{4}.
\end{equation}
Indeed their generating function is
$\sum_{a,b}w_{ab}z^{a+b}=S(z)/z$.
If just one low count $m_a$ is differentiated $s\geq1$ times, then
\[
 \chi_{\mathbf r}=2M_{a,r_D}
        +(2^s-2)w_{aa}(2a)^{r_D}.
\]
If two distinct low counts $m_a,m_b$ are differentiated positively,
the answer is $2w_{ab}(a+b)^{r_D}$; with three distinct low counts it
is zero. Powers with exponent zero are one, including $0^0$ in these
finite polynomial evaluations.

For explicit reproduction, write $X=(p,D,m_0,m_1,m_2)$ and
$N(\mathbf r)=\mathcal L[X^{\mathbf r}]$. Let $\mathbf e_i$ be
the $i$th unit vector for $1\leq i\leq5$. Starting with $N(0)=1$,
the usual product rule gives the finite recurrence
\begin{equation}\label{eq:a:bell-rule}
 N(\mathbf r+\mathbf e_i)=
 \sum_{\mathbf s\leq\mathbf r}\binom{\mathbf r}{\mathbf s}
       \chi_{\mathbf s+\mathbf e_i}N(\mathbf r-\mathbf s).
\end{equation}
For example $\mathcal L[D]=-S(1)$ and
$\mathcal L[D^2]=S(1)^2+S(1)$. Also
$\mathcal L[m_0]=2M_{0,0}=8$; its second logarithmic derivative is
$8+2w_{00}=8+2\pi^2$, whence
$\mathcal L[m_0^2]=72+2\pi^2$.
As a further useful check,
\[
 \mathcal L[D+p]=\mathcal L[(D+p)^2]=0,\qquad
 \mathcal L[(D+p)^3]=\frac{\pi^2}{2},\quad
 \mathcal L[(D+p)^4]=\frac{3\pi^2}{4}.
\]

We now apply these rules to the pure-loop core. In this and the next
calculation, a vertical restriction $\big|_{q=r}$ means the part of
the exact graph sum whose distinguished large vertex has core
complexity $r$, as defined in Section~\ref{subsec:s03:04}.
For the pure-loop core its restored vertex has $(G,n)=(g-p,2p)$ and
complexity $2g-2$. Put, locally,
\[
 E_r=e_{2p,r}(a_1,b_1,\ldots,a_p,b_p)\quad(r=1,2),
\]
\[
 v_1=\frac p2+c_1(2p),\qquad
 v_2=\frac{3p^2}{8}+\frac{3p}{2}c_1(2p)+c_2(2p).
\]
The square-root genus factor and the absolute fixed-marking volume
expansion give these $v_1,v_2$. Multiplication by the profile in
complexity $2g-2$ gives respectively
\begin{equation}\label{eq:a:pure-loop-polynomials}
 v_1+\frac{E_1}{2},\qquad
 v_2+\frac{(1+v_1)E_1}{2}+\frac{E_2}{4}.
\end{equation}
The additional $E_1/2$ in the second expression comes from expanding
$(2g-2)^{-1}$; omitting it changes the second coefficient.

\begin{prp}[Pure-loop coefficients]\label{prp:a:pure-loops}
The pure-loop contribution in the absolute normalization is
\begin{equation}\label{eq:a:q-zero}
\begin{aligned}
 \left.\frac{Z_g(1)}{F_g}\right|_{q=0}
 &=\frac{e^{S(1)}}{\sqrt\pi}\Bigl[1
   +\frac{1/12-65/(6\pi^2)}g\\
 &\quad+\frac{43/288-469/(36\pi^2)+5665/(72\pi^4)}{g^2}
   +O(g^{-3})\Bigr].
\end{aligned}
\end{equation}
\end{prp}
\begin{prf}
Sum~\eqref{eq:a:pure-loop-polynomials} with $\mathcal L$.
For the first coefficient the finite rules give
\[
 \mathcal L[v_1]=c_1(0)+\frac{3S(1)-2S(1)^2}{\pi^2},\qquad
 \mathcal L[E_1]=-1+\frac{4S(1)^2-6S(1)-16}{\pi^2}.
\]
Here is a compact coefficient table for the second calculation. Each
column lists the coefficient of $S(1)^j$ in the indicated expression,
after multiplying the entire expression by $96\pi^4$:
\[
\begin{array}{c|r|r}
 j&\mathcal L[v_2-c_2(0)]&\mathcal L[(1+v_1)E_1]\\\hline
0&0&-152\pi^4-2448\pi^2+2816\\
1&-3\pi^4+568\pi^2-2064&-816\pi^2+13728\\
2&-400\pi^2+1072&800\pi^2+3328\\
3&-16\pi^2-384&-768\\
4&192&-768
\end{array}
\]
\[
\begin{array}{c|r}
 j&\mathcal L[E_2]\\\hline
0&168\pi^4+2000\pi^2+16512\\
1&12\pi^4-640\pi^2-19200\\
2&-10944\\
3&64\pi^2+3072\\
4&768
\end{array}
\]
The table follows by inserting~\eqref{eq:a:first-profile},
\eqref{eq:a:second-profile}, and~\eqref{eq:a:c-two-n} into
\eqref{eq:a:bell-rule}; the largest polynomial degree is four.
Combining its columns with coefficients $1,1/2,1/4$ cancels every
positive power of $S(1)$. The constant left is
\[
 c_2(0)-\frac{17}{48}-\frac{181}{24\pi^2}
                                      +\frac{173}{3\pi^4}.
\]
Proposition~\ref{prp:a:c-two} now gives the displayed result.
Section~\ref{sec:s03} justifies the operation on all legal loop
counts: its degree-moment estimate and summable core remainder apply
to these fixed-degree polynomials, and the large-loop tail is smaller
than every fixed inverse power of $g$. Thus extending the finite legal
sum to the signed entire series changes only the stated remainder.
\end{prf}

For complexity one there are three cores: a rational vertex with a
small loop and one connector to the root, a rational vertex with three
connectors, and an elliptic vertex with one connector. Their precise
graphs and automorphisms will be listed in
Subsection~\ref{subsec:a:04}. We calculate their first corrections
here so that the finite core table has a proved coefficient rule.

For a restored core with root valence $n_0$, root genus loss
$c=(n_0+1)/2$, and $p$ restored root loops, the absolute volume
correction is
\[
 \frac32+\frac{p+c}{2}+c_1(2p+n_0).
\]
The $3/2$ is the first correction of the complexity-one factorial
ratio. Let $D_{\rm c}$ be the total degree at its connector flags on
the root, and $m_{0,\rm c}$ the number of their zero degrees. Add them
to the loop statistics in the first profile. The same finite moment
rule gives, for $n_0=1$ or $3$,
\begin{align}
 &\mathcal L\left[\frac32+\frac{p+c}{2}+c_1(2p+n_0)
       +\frac12e_{2p+n_0,1}(\text{loop and connector degrees})\right]
 \notag\\
 &\quad=\frac{22\pi^2-12D_{\rm c}^2+(30-12n_0)D_{\rm c}
      -3m_{0,\rm c}^2-33m_{0,\rm c}+3n_0-121}{12\pi^2}.
 \label{eq:a:q-one-loop-average}
\end{align}
Here the profile's total degree is $D+D_{\rm c}$ and its zero count
is $m_0+m_{0,\rm c}$, with $n=2p+n_0$.

Only elementary connector moments remain. For one connector let its
degree be $a$ and its zero indicator be $z=\mathbf1_{a=0}$. Sum its
negative-psi coefficient against $1$ or
$W^{\rm WP}_{1,1}(t)=(2\pi^2-t)/24$, respectively. Use $x$ locally
as a complex degree-generating parameter near $1$. For degree orders
$r=0,1,2$, the raw moments are the unscaled integrals
\[
\begin{aligned}
 \text{rational:}\quad &2M_{0,r}
 =\left.(x\partial_x)^r\int_0^\alpha f(xt)\,\dd t\right|_{x=1},\\
 \text{elliptic:}\quad &\left.(x\partial_x)^r
   \int_0^\alpha W^{\rm WP}_{1,1}(t)f(xt)\,\dd t\right|_{x=1}.
\end{aligned}
\]
The elliptic polynomial $W^{\rm WP}_{1,1}(t)$ is not scaled with $x$. A connector
has no loop branch-exchange factor $1/2$. The zero-indicator moments
are separately
\[
 \int_0^\alpha 1\,\dd t=2\pi^2,\qquad
 \int_0^\alpha W^{\rm WP}_{1,1}(t)\,\dd t=\frac{\pi^4}{12}.
\]
These give the following raw moment table:
\begin{equation}\label{eq:a:connector-moments}
\begin{array}{c|rrrr}
 &1&a&a^2&z\\\hline
 \text{rational connector}&8&-8&8-\pi^2&2\pi^2\\
 \text{elliptic connector}&(\pi^2+4)/3&-8/3&(16-2\pi^2)/3&\pi^4/12
\end{array}
\end{equation}
Always $z^r=z$
for $r\geq1$, and $a^r z=0$ for $r\geq1$. For $m$ independent
connectors with raw moments $h_0,h_1,h_2,h_z$, product expansion gives
\begin{align*}
 \sum D_{\rm c}&=mh_1h_0^{m-1},&
 \sum D_{\rm c}^2&=mh_2h_0^{m-1}+m(m-1)h_1^2h_0^{m-2},\\
 \sum m_{0,\rm c}&=mh_zh_0^{m-1},&
 \sum m_{0,\rm c}^2&=mh_zh_0^{m-1}+m(m-1)h_z^2h_0^{m-2}.
\end{align*}
For $m=1$ the terms with $m(m-1)$ are absent. These sums mean the
product of the unnormalized connector weights, not a probability
average.

\begin{prp}[Complexity-one coefficients]\label{prp:a:q-one}
After removing the common factor $e^{S(1)}/\sqrt\pi$, the three
complexity-one cores contribute the following coefficients of
$g^{-1}$ and $g^{-2}$ in $Z_g(1)/F_g$:
\begin{equation}\label{eq:a:q-one-table}
\begin{array}{c|c|c}
 \text{Small vertex}&g^{-1}&g^{-2}\\\hline
 \begin{gathered}(0,3)\\\text{one connector}\\\text{and one loop}\end{gathered}
 &1&29/24-37/(3\pi^2)\\
 \begin{gathered}(0,3)\\\text{three connectors}\end{gathered}
 &32/(3\pi^2)&-1-4/(9\pi^2)-1616/(9\pi^4)\\
 \begin{gathered}(1,1)\\\text{one connector}\end{gathered}
 &1/24+1/(6\pi^2)&13/288-1/(48\pi^2)-101/(36\pi^4)
\end{array}
\end{equation}
Consequently the complete complexity-one part is
\begin{equation}\label{eq:a:q-one}
\begin{aligned}
 \left.\frac{Z_g(1)}{F_g}\right|_{q=1}
 &=\frac{e^{S(1)}}{\sqrt\pi}\Bigl[
   \frac{25/24+65/(6\pi^2)}g\\
 &\quad+\frac{73/288-1843/(144\pi^2)-6565/(36\pi^4)}{g^2}
   +O(g^{-3})\Bigr].
\end{aligned}
\end{equation}
\end{prp}
\begin{prf}
Insert the connector moments into~\eqref{eq:a:q-one-loop-average}.
For the three rows in order use $(n_0,m)=(1,1),(3,3),(1,1)$.
The remaining prefactors are $1/8$, $1/(48\pi^2)$, and
$1/(8\pi^2)$, respectively. The first includes the small-loop integral
$\alpha$ and its automorphism factor $2$; the second includes $3!$;
the third has no graph symmetry. Multiplying the constant connector
moment gives the leading column, and multiplying the polynomial in
\eqref{eq:a:q-one-loop-average} gives the correction column.
For example, the first polynomial at $n_0=1$ has raw average
\[
 \frac{8(22\pi^2-118)-12(8-\pi^2)+18(-8)
                   -36(2\pi^2)}{12\pi^2}
 =\frac{29}{3}-\frac{296}{3\pi^2}.
\]
Its factor $1/8$ gives the first correction in the table. The two
other rows follow by the displayed product formulas, with no
uncomputed endpoint moments. Summing proves~\eqref{eq:a:q-one}.
The summable remainders are again those of
\eqref{eq:s03:summable-remainder}, now to first order on a fixed
complexity-one core.
\end{prf}

\begin{rmk}[Normalization]
The correction columns are absolute coefficients after extracting
$e^{S(1)}/\sqrt\pi$ from $Z_g(1)/F_g$. The functions $\gamma_r$ of
Section~\ref{sec:s03} are instead coefficients of $Z_g/V_{g,0}$.
If the complete absolute coefficients are $r_1,r_2$ (in this remark
only), then
\[
 e^{-S(1)}\gamma_1(1)=r_1-c_1(0),\qquad
 e^{-S(1)}\gamma_2(1)=r_2-c_2(0)-c_1(0)(r_1-c_1(0)).
\]
Thus an absolute core correction must not be relabelled as a relative
$\gamma_r$ coefficient.
\end{rmk}
